% !TeX root = ../sections/02_exercises.tex
\subsection{2-B-5.}
\begin{exercise}
	数列 \(\{a_n\}_{n\in\mathbb{N}}\) が有界でないならば，
	\[
	\lim_{k\to\infty}|a_{n_k}|=\infty
	\]
	となる部分列
	\[
	\{a_{n_k}\}_{k\in\mathbb{N}}
	\]
	が存在することを証明せよ。
\end{exercise}

\begin{background}
	数列 \(\{a_n\}\) が有界であるとは，
	ある \(M>0\) が存在して，任意の \(n\in\mathbb{N}\) に対して
	\[
	|a_n|\leq M
	\]
	が成り立つことをいう。
	
	したがって，\(\{a_n\}\) が有界でないとは，
	任意の \(M>0\) に対して，ある \(n\in\mathbb{N}\) が存在して
	\[
	|a_n|>M
	\]
	となることを意味する。
	
	ただし，部分列を作るためには，添字が
	\[
	n_1<n_2<n_3<\cdots
	\]
	となるように選ぶ必要がある。
	そのためには，「どの番号以降にも大きな項が存在する」
	ことを確認してから，順に \(n_k\) を選ぶ。
\end{background}

\begin{strategy}
	まず，任意の \(N\in\mathbb{N}\) に対して，尾部
	\[
	\{a_n\mid n\geq N\}
	\]
	が有界でないことを示す。
	
	もしある \(N\) 以降の尾部が有界ならば，
	ある \(M_1>0\) が存在して
	\[
	n\geq N
	\quad\Longrightarrow\quad
	|a_n|\leq M_1
	\]
	となる。
	また，有限個の項
	\[
	a_1,a_2,\ldots,a_{N-1}
	\]
	も有界である。
	したがって，数列全体が有界になってしまい，仮定に反する。
	
	よって，任意の \(N\) 以降にはいくらでも大きな項がある。
	
	これを用いて，帰納的に部分列を作る。
	まず，尾部全体が有界でないので，
	ある \(n_1\) が存在して
	\[
	|a_{n_1}|>1
	\]
	となる。
	
	次に，\(n_1+1\) 以降の尾部も有界でないので，
	ある \(n_2>n_1\) が存在して
	\[
	|a_{n_2}|>2
	\]
	となる。
	
	同様にして，すでに
	\[
	n_1<n_2<\cdots<n_{k-1}
	\]
	を選んだとき，\(n_{k-1}+1\) 以降の尾部が有界でないことから，
	ある \(n_k>n_{k-1}\) が存在して
	\[
	|a_{n_k}|>k
	\]
	となるように選べる。
	
	このようにして作った部分列は，任意の \(k\) に対して
	\[
	|a_{n_k}|>k
	\]
	を満たす。
	したがって，
	\[
	k\to\infty
	\quad\Longrightarrow\quad
	|a_{n_k}|\to\infty
	\]
	である。
	
	よって，
	\[
	\lim_{k\to\infty}|a_{n_k}|=\infty
	\]
	となる部分列が存在する。
\end{strategy}